\documentclass[10pt,reqno]{amsart}

\usepackage[utf8]{inputenc}
\usepackage[T1]{fontenc}
\usepackage{lmodern}
\usepackage{amsmath,amssymb,amsthm,mathtools}
\usepackage{array,booktabs}
\usepackage[margin=0.82in]{geometry}
\usepackage{xcolor}
\usepackage[colorlinks=true,linkcolor=blue!50!black,citecolor=blue!50!black]{hyperref}

\newtheorem{theorem}{Theorem}[section]
\newtheorem{proposition}[theorem]{Proposition}
\newtheorem{lemma}[theorem]{Lemma}
\theoremstyle{definition}
\newtheorem{definition}[theorem]{Definition}
\newtheorem{remark}[theorem]{Remark}

\newcommand{\chiad}{\chi_{\mathrm{ad}}}
\newcommand{\Qcone}{\mathcal Q}
\newcolumntype{L}[1]{>{\raggedright\arraybackslash}p{#1}}

\title[Adjoint-generated Ginibre $Q_3$: Part III]
{The full adjoint-generated case of Ginibre's condition
\texorpdfstring{$Q_3$}{Q3}\\
\large Part III: signed-product completion}
\author{Leslie P. Polzer}
\address{}
\email{polzer@fastmail.com}
\date{July 2026}

\hypersetup{
  pdftitle={The full adjoint-generated case of Ginibre's condition Q3: Part III},
  pdfauthor={Leslie P. Polzer},
  pdfsubject={Reader proof spine for the full signed-product hierarchy},
  pdfkeywords={Ginibre Q3, compact Lie group, adjoint character, exact certificates}
}

\begin{document}

\begin{abstract}
Let $G$ be a compact connected Lie group with simple Lie algebra and let
$\Qcone_G$ be the closed multiplicative convex cone generated by its real
adjoint character.  We reduce Ginibre's full condition $Q_3$ on $\Qcone_G$
to an explicit two-parameter moment hierarchy and close that hierarchy for
every simple type.  Parts~I--II prove the two-minus row.  The present reader
part states the exact higher-minus reduction, the bounded-Littlewood moment
contract, and the exhaustive classification certificate.  Detailed tail
estimates, finite ledgers, verifier semantics, and replay instructions are in
the separately authenticated 35-page formal and computational supplement
\texttt{full\_q3\_extension.pdf}; that supplement is part of the proof, not
an optional derivation archive.
\end{abstract}

\maketitle

\begin{remark}[Reader manuscript and formal supplement]
\label{rem:architecture}
The reader manuscript \path{submission.pdf} consists of Parts~I--II
(\path{paper.pdf}) and this compact Part~III
(\path{full_q3_main.pdf}).  The proof package also contains the formal and
computational supplement \path{full_q3_extension.pdf}.  The artifact manifest
\path{publication_artifacts.sha256} binds all three components and the reader
manuscript.  References below to ``the supplement'' mean that exact file.
The optional \path{paper_full.pdf} development archive is not load-bearing.
\end{remark}

\section{Reduction to the higher-minus hierarchy}

Let $\mu_G$ be normalized Haar measure, and let
$X=\chiad(g)$ and $Y=\chiad(h)$ for independent Haar variables.

\begin{definition}
Set
\[
 \Qcone_G=\overline{\left\{\sum_{j=0}^{d}c_j\chiad^j:
 d\geq0,\ c_j\geq0\right\}}^{\|\cdot\|_\infty},
 \qquad
 I^G_{a,n}=\mathbf E\,(X-Y)^{2a}(X+Y)^n.
\]
Every member of $\Qcone_G$ is a real continuous positive-definite function:
characters are positive definite, and products, nonnegative sums, and
uniform limits preserve that property.
\end{definition}

\begin{proposition}[Exact reduction and cone promotion]
\label{prop:reduction}
Ginibre's condition $Q_3$ holds on $\Qcone_G$ if and only if
$I^G_{a,n}\geq0$ for all $a,n\geq0$.  Only $a\geq1$ and odd $n$ are
nonautomatic, and Parts~I--II prove the complete row $a=1$.
\end{proposition}

\begin{proof}
A signed product of the single generator, with $r$ minus factors and $n$
plus factors, is $\mathbf E(X-Y)^r(X+Y)^n$.  It vanishes for odd $r$ by
interchanging $X,Y$, and equals $I^G_{a,n}$ for $r=2a$.  Even $n$ gives a
pointwise nonnegative integrand.  For $a=0$, expansion gives a sum of
nonnegative invariant-tensor multiplicities
$m_j^G=\dim((\mathfrak g_{\mathbb C}^{\otimes j})^G)$.

The promotion from $\chiad$ to its generated cone uses
\[
 (fg)(x)\mathbin\pm(fg)(y)=\tfrac12\{(f(x)+f(y))(g(x)\mathbin\pm g(y))
 +(f(x)-f(y))(g(x)\mathbin\mp g(y))\}.
\]
Iteration, nonnegative linearity, and uniform closure give precisely
Ginibre's Proposition~1 \cite{Ginibre1970}.  Finally, the $a=1$ assertion is
the main theorem of Parts~I--II \cite{AdjointTwoMinus}.
\end{proof}

\begin{lemma}[Exact moment functional]
\label{lem:moment}
For $m_j^G=\mathbf E X^j$,
\[
 I^G_{a,n}=\sum_{j=0}^{2a}\sum_{k=0}^{n}(-1)^j
 \binom{2a}{j}\binom nk,m^G_{2a-j+k}m^G_{j+n-k}.
\]
Equivalently, if $M_G(t)=\mathbf E e^{tX}$, then
\[
 \sum_{n\geq0}I^G_{a,n}\frac{t^n}{n!}
 =\sum_{j=0}^{2a}(-1)^j\binom{2a}{j}
 M_G^{(2a-j)}(t)M_G^{(j)}(t).
\]
\end{lemma}

\begin{proof}
Expand the two binomials.  The generating-function identity follows by
expanding only $(X-Y)^{2a}$ after inserting $e^{t(X+Y)}$.
\end{proof}

\begin{proposition}[Why the cone cannot be enlarged indiscriminately]
For $G=\mathrm{SO}(3)$, the cone of all real continuous positive-definite
functions does not satisfy $Q_3$.  There are four matrix coefficients
$f_v(g)=\langle v,gv\rangle$ and the sign pattern $--++$ for which the
integral is $-8/279$.
\end{proposition}

\begin{proof}
Each $f_v$ is positive definite since its defining Gram quadratic form is
$\|\sum_j z_jg_jv\|^2$.  Haar projection on
$(\mathbb R^3)^{\otimes2}$ and $(\mathbb R^3)^{\otimes4}$ gives the claimed
rational value for the explicit four vectors printed in Supplement~\S1.
Thus this proposition concerns a genuinely larger cone than $\Qcone_G$.
\end{proof}

\section{Exact finite-rank contract}

The higher-minus proof partitions each simple family into a stable prefix,
an exact finite box, and a proved tail.  The only substantial finite-rank
moment interface not already printed in Parts~I--II is the following
bounded-Littlewood formula.

\begin{lemma}[Bounded-Littlewood moment determinants]
\label{lem:bl}
In the completed symmetric-function ring, put $e_r=h_r=0$ for $r<0$ and
\[
 f_r^e=\sum_{k\in\mathbb Z}e_ke_{k+r},\quad
 f_r^h=\sum_{k\in\mathbb Z}h_kh_{k+r},\quad
 e=\sum_{k\geq0}e_k,\quad \bar e=\sum_{k\geq0}(-1)^ke_k.
\]
For $b\geq1$, define
\[
\begin{aligned}
 \mathcal B_b&=\tfrac12\{e\det(f^e_{i-j}-f^e_{i+j-1})
 +\bar e\det(f^e_{i-j}+f^e_{i+j-1})\},\\
 \mathcal C_b&=\det(f^h_{i-j}-f^h_{i+j}),\qquad
 \mathcal D_b=\tfrac12\det(f^e_{i-j}+f^e_{i+j-2}),
\end{aligned}
\]
where all determinants are $b\times b$.  Then
\[
 m_j^{B_b}=[m_{(2^j)}]\mathcal B_b,\qquad
 m_j^{C_b}=[m_{(2^j)}]\mathcal C_b,\qquad
 m_j^{D_b}=[m_{(2^j)}]\mathcal D_b.
\]
If $F$ is homogeneous of degree $2j$ and
$\pi_2(p_1)=p$, $\pi_2(p_2)=q$, $\pi_2(p_r)=0$ for $r\geq3$, then
\[
 [m_{(2^j)}]F=j!\sum_{a=0}^{j}(2(j-a)-1)!!
 [p^{2(j-a)}q^a]\pi_2(F).
\]
\end{lemma}

\begin{proof}
The orthogonal and symplectic Schur-integral criteria identify the admitted
even-column/even-row sectors \cite{Rains1997}.  The bounded Littlewood
identities of Huh--Kim--Krattenthaler--Okada give the three displayed
determinants \cite{HKKO2025}; Hall duality turns their Schur sums into the
stated adjoint moments.  In the last formula, only
$p_1^{2(j-a)}p_2^a$ can contribute to $m_{(2^j)}$; counting assignments of
the ordered factors gives $j!(2(j-a)-1)!!$.  Supplement~\S8 supplies the
degree-by-degree derivation, specializations, and recurrences used by the
exact verifier.
\end{proof}

\begin{proposition}[Exhaustive certificate contract]
\label{prop:contract}
For every simple type, every $a\geq2$, and every odd $n\geq1$, one has
$I^G_{a,n}>0$.  The disjoint stable, finite, and tail regions are those in
the following audit summary; the supplement states every endpoint and proves
that the regions cover without gaps.
\[
\begin{array}{@{}L{0.19\textwidth}L{0.39\textwidth}L{0.31\textwidth}@{}}
\toprule
\text{family}&\text{exact finite supplier}&\text{stable/tail supplier}\\
\midrule
$A_1$&$A_1$: 78 values; minimum $16$&effective Haar rectangle\\
$A_2$&$A_2$: 78 values; minimum $72$&exact moment cap\\
$A_3,A_4$&66 and 136 values; minima $94,96$&exact moment caps\\
$A_r$, $r\geq5$&1,315 nonstable values; minimum $3,550$&stable Schur--Weyl row and uniform tail\\
$G_2,F_4,E_6,E_7,E_8$&1,265 values; minimum $36$&five exact cap tails\\
$B_b,C_b,D_b$&17,862 values; minimum $I^{D_5}_{2,1}=50$&500 stable identities, exact middle/high and 58 directed low-row tails\\
\bottomrule
\end{array}
\]
The classical finite supplier covers exactly
\[
 B_2,\ldots,B_{21},\qquad C_2,\ldots,C_{28},\qquad
 D_4,\ldots,D_{70}.
\]
It reconstructs moments through degree $71$ by Lemma~\ref{lem:bl}, checks
all $17{,}862$ values with GMP integers, and uses a 961-bit CRT modulus.
An independent checker agrees on the unresolved high-degree moments modulo
21 primes (628 bits, exceeding the 607-bit uniqueness bound) and recomputes
the remaining 5,509 inequalities.  The least directed low-tail logarithmic
margin is $0.0971194051797054477804174\ldots$ at $C_{15}$.
\end{proposition}

\begin{proof}
This is the compact statement of Supplement~\S\S2--8.  Its type-$A$ proofs
combine Schur--Weyl moment formulas, exact cap inequalities, and the stated
finite GMP boxes.  Its exceptional proof uses exact invariant-tensor
moments, five cap tails, and the five exhaustive residual boxes.  For the
classical families, stable cumulants give the prefix; analytic trace and Weyl
cap bounds give the middle and high tails; directed 384-bit interval
arithmetic proves all 58 remaining row tails; and Lemma~\ref{lem:bl} supplies
the exact finite moments.  The supplement prints the row onsets and proves
the three regions exhaustive.  Each quoted count, minimum, margin, and scope
is asserted by its source verifier and authenticated transcript.  Therefore
the proposition is a finite consequence of the displayed mathematical
interfaces and the supplement's exact certificate ledger.
\end{proof}

\section{Full theorem}

\begin{theorem}[Full adjoint-generated Ginibre $Q_3$]
\label{thm:main}
Let $G$ be any compact connected Lie group with simple Lie algebra.  For
every finite tuple $f_1,\ldots,f_r\in\Qcone_G$ and signs
$\varepsilon_i\in\{+1,-1\}$ containing an even number of minus signs,
\[
 \int_{G\times G}\prod_{i=1}^{r}
 \bigl(f_i(g)+\varepsilon_i f_i(h)\bigr)
 \,d\mu_G(g)d\mu_G(h)\geq0.
\]
If the number of minus signs is odd, the integral is zero.  Thus
$\Qcone_G$ satisfies Ginibre's condition $Q_3$ with its full finite-tuple
and sign-pattern quantifiers.
\end{theorem}

\begin{proof}
Proposition~\ref{prop:reduction} supplies the automatic sectors, imports the
two-minus row, and promotes the generator hierarchy to $\Qcone_G$.
Proposition~\ref{prop:contract} supplies every remaining $a\geq2$, odd-$n$
case in types $A,B,C,D$ and the five exceptional types.  Classification and
central-cover invariance exhaust all compact connected global forms.
\end{proof}

\begin{remark}[Reproduction]
From the source root, \texttt{make publication-artifact-audit} rebuilds and
hash-checks the reader manuscript and supplement;
\texttt{make full-q3-extension} replays the complete mathematical
certificate suite.  The CI workflow runs the publication build, the clean
room Parts~I--II replay, and the independent Part~III replay as separate
jobs.  Exact commands, trusted-base boundaries, and transcript mappings are
in \path{REPLAY.md} and Supplement~\S9.
\end{remark}

\begin{thebibliography}{9}

\bibitem{AdjointTwoMinus}
L.~P. Polzer,
The two-minus adjoint-character case of Ginibre's $Q_3$ condition,
Parts I--II: classification and exact certificates, 2026.

\bibitem{Ginibre1970}
J.~Ginibre,
General formulation of Griffiths' inequalities,
\emph{Comm. Math. Phys.} \textbf{16} (1970), 310--328.

\bibitem{HKKO2025}
J.~Huh, J.~S. Kim, C.~Krattenthaler, and S.~Okada,
Bounded Littlewood identities with fixed number of odd rows or odd columns,
\emph{Electron. J. Combin.} \textbf{33} (2026), no.~3, Paper~P3.5.

\bibitem{Rains1997}
E.~M. Rains,
Increasing subsequences and the classical groups,
\emph{Electron. J. Combin.} \textbf{5} (1998), Research Paper~12.

\end{thebibliography}

\end{document}
